\documentclass[11pt,a4paper]{article}

% ---------------------------------------------------------------------------
% Relatorio PTOSS-2 - Testes Caixa-Branca do fluxograma_controller
% NoFluxoUNB - FGA0314 Testes de Software
% Compilar com: pdflatex relatorio-ptoss2-caixa-branca.tex
% ---------------------------------------------------------------------------

\usepackage[utf8]{inputenc}
\usepackage[T1]{fontenc}
\usepackage[brazil]{babel}
\usepackage[margin=2.5cm]{geometry}
\usepackage{booktabs}
\usepackage{array}
\usepackage{xcolor}
\usepackage{listings}
\usepackage{hyperref}
\usepackage{enumitem}
\usepackage{parskip}

\hypersetup{colorlinks=true, linkcolor=violet, urlcolor=violet}

\definecolor{codebg}{RGB}{245,245,248}
\definecolor{codecomment}{RGB}{110,110,120}
\lstset{
  backgroundcolor=\color{codebg},
  basicstyle=\ttfamily\small,
  commentstyle=\color{codecomment},
  breaklines=true, frame=single, rulecolor=\color{lightgray},
  framesep=6pt, xleftmargin=6pt, showstringspaces=false,
  columns=fullflexible, keepspaces=true
}

\newcommand{\V}{\textbf{V}}
\newcommand{\Ff}{\textbf{F}}

\title{\textbf{Relatório PTOSS-2 --- Testes de Caixa-Branca\\ do \texttt{fluxograma\_controller}}}
\author{NoFluxoUNB \\ \small FGA0314 --- Testes de Software (2026.1)}
\date{\today}

\begin{document}
\maketitle

\begin{abstract}
\noindent
Este relatório descreve a aplicação da técnica de \textbf{teste estrutural
(caixa-branca)} sobre as funções auxiliares do \texttt{fluxograma\_controller},
o módulo mais complexo do backend do NoFluxoUNB. Apresentamos a motivação, o
projeto dos casos de teste no formato de três tabelas (Condições,
Tabela-Verdade e Casos) com o critério \textbf{MC/DC}, e os resultados de
cobertura obtidos. Esta entrega corresponde à \textbf{Fase 2} do nosso plano de
testes (aumento de cobertura).
\end{abstract}

\tableofcontents
\bigskip

% ===========================================================================
\section{Por que fizemos esses testes (e como ajudam o projeto)}
% ===========================================================================

O \texttt{fluxograma\_controller} ($\approx$1.000 linhas) é o coração do produto:
é ele que \emph{casa} as disciplinas do histórico do aluno com as matérias do
curso, resolve \emph{equivalências} e \emph{pré-requisitos} e calcula o que já
foi integralizado. Se essa lógica erra, o aluno vê um fluxograma errado --- uma
matéria aparece como concluída sem ter sido, ou como pendente quando já foi
cumprida por equivalência. É exatamente o tipo de erro que mais prejudica a
confiança no sistema.

Por isso escolhemos a técnica de \textbf{caixa-branca}: diferente da caixa-preta
(que olha só a especificação), aqui exercitamos os \textbf{caminhos internos do
código} --- cada decisão, cada ramo, cada condição de uma expressão composta.
Isso encontra classes de defeito que a caixa-preta dificilmente pega, como:

\begin{itemize}[leftmargin=1.4em]
  \item erros em \textbf{condições compostas} e \textbf{curto-circuito}
        (ex.: \texttt{tipo\_natureza !== undefined \&\& !== null});
  \item ramos de \textbf{fallback} raramente percorridos
        (ex.: os 4 níveis de \texttt{findSubjectMatch});
  \item lógica de \textbf{prioridade} na resolução de duplicatas
        (\texttt{processMatchedDiscipline}).
\end{itemize}

\paragraph{Benefícios concretos para o projeto.}
\begin{enumerate}[leftmargin=1.6em]
  \item \textbf{Rede de segurança contra regressão}: os testes rodam no CI a
        cada PR; se alguém quebrar a lógica de casamento/equivalência, o build
        falha antes do merge.
  \item \textbf{Documentação executável}: cada caso registra, de forma
        verificável, qual é o comportamento esperado de cada função.
  \item \textbf{Testes rápidos e determinísticos}: testam funções puras, sem
        banco de dados nem rede --- ao contrário dos testes de endpoint, que
        dependem de \emph{mocks} frágeis do Supabase.
  \item \textbf{Cobertura mensurável}: elevamos a cobertura do módulo de
        $\sim$10\% para $\sim$28\% (ver Seção~\ref{sec:resultados}), cumprindo a
        Fase 2 do plano de testes.
\end{enumerate}

% ===========================================================================
\section{Alvo dos testes}
% ===========================================================================

As funções auxiliares são internas ao módulo. Para testá-las de forma isolada
(sem passar pelos \emph{endpoints} HTTP), nós as expusemos por meio de um export
marcado, usado \textbf{apenas} pelos testes, sem alterar o comportamento de
produção:

\begin{lstlisting}[language=Java]
// fluxograma_controller.ts
export const __testing__ = {
    isOptativa, mapPreRequisitosFromDb, extractSubjectCodes,
    getCodigosEquivalentes, codigoContidoNaEquivalencia, getStatusPriority,
    findSubjectMatch, processMatchedDiscipline, checkEquivalencies,
    mapEquivalenciasFromDb,
};
\end{lstlisting}

% ===========================================================================
\section{Exemplo completo de MC/DC: \texttt{isOptativa}}
% ===========================================================================

Uma matéria é optativa se \texttt{tipo\_natureza == 1}; quando esse campo não
está definido, usa-se o \emph{fallback} \texttt{nivel == 0}. O código é:

\begin{lstlisting}[language=Java]
function isOptativa(mpc): boolean {
    if (mpc.tipo_natureza !== undefined && mpc.tipo_natureza !== null)
        return mpc.tipo_natureza === 1;
    return (mpc.nivel ?? 0) === 0;
}
\end{lstlisting}

A decisão principal é $D = C_1 \land C_2$. Aplicamos MC/DC nas três tabelas no
formato da disciplina.

\subsection*{Tabela 1 --- Condições}

\begin{table}[h]
\centering
\small
\begin{tabular}{@{}cll@{}}
\toprule
\textbf{Id} & \textbf{Condição / Decisão} & \textbf{Descrição} \\
\midrule
$D$   & \texttt{D = C1 \&\& C2}            & decide se usa \texttt{tipo\_natureza} ou o fallback \\
$C_1$ & \texttt{tipo\_natureza !== undefined} & o campo foi informado \\
$C_2$ & \texttt{tipo\_natureza !== null}      & o campo não é nulo \\
\bottomrule
\end{tabular}
\caption{Decisão e condições atômicas de \texttt{isOptativa}.}
\end{table}

\subsection*{Tabela 2 --- Tabela-Verdade e pares de independência}

\begin{table}[h]
\centering
\small
\begin{tabular}{@{}ccccl@{}}
\toprule
\textbf{\#} & $C_1$ & $C_2$ & $D$ & \textbf{Observação} \\
\midrule
1 & \V  & \V  & \V  & \texttt{tipo\_natureza} definido e não nulo \\
2 & \V  & \Ff & \Ff & \texttt{tipo\_natureza = null} \\
3 & \Ff & \V  & \Ff & \texttt{tipo\_natureza = undefined} (curto-circuito) \\
4 & \Ff & \Ff & \Ff & combinação impossível na prática \\
\bottomrule
\end{tabular}
\caption{Tabela-verdade de $D = C_1 \land C_2$. Pares de independência:
\textbf{(1,3)} isola $C_1$ ($C_2$ fixo em \V, $D$ muda) e
\textbf{(1,2)} isola $C_2$ ($C_1$ fixo em \V, $D$ muda).}
\end{table}

O conjunto MC/DC mínimo que cobre os dois pares é $\{1, 2, 3\}$. Cada linha gera
ainda uma subdecisão (\texttt{tipo\_natureza === 1} quando $D=\V$;
\texttt{(nivel ?? 0) === 0} quando $D=\Ff$), que também exercitamos.

\subsection*{Tabela 3 --- Casos de Teste}

\begin{table}[h]
\centering
\small
\begin{tabular}{@{}clll@{}}
\toprule
\textbf{CT} & \textbf{Entrada} & \textbf{Cobre} & \textbf{Saída esperada} \\
\midrule
CT1 & \texttt{\{tipo\_natureza: 1\}}                & linha 1 ($D$\V), sub \V & \texttt{true} (optativa) \\
CT2 & \texttt{\{tipo\_natureza: 0\}}                & linha 1 ($D$\V), sub \Ff & \texttt{false} (obrigatória) \\
CT3 & \texttt{\{tipo\_natureza: null, nivel: 0\}}   & linha 2 ($D$\Ff), fb \V  & \texttt{true} \\
CT4 & \texttt{\{tipo\_natureza: null, nivel: 5\}}   & linha 2 ($D$\Ff), fb \Ff & \texttt{false} \\
CT5 & \texttt{\{tipo\_natureza: undefined, nivel: 0\}} & linha 3 ($D$\Ff), fb \V & \texttt{true} \\
CT6 & \texttt{\{tipo\_natureza: undefined, nivel: 5\}} & linha 3 ($D$\Ff), fb \Ff & \texttt{false} \\
CT7 & \texttt{\{\}} (sem \texttt{nivel})            & \texttt{nivel ?? 0 = 0}  & \texttt{true} \\
\bottomrule
\end{tabular}
\caption{Casos de teste de \texttt{isOptativa} (fb = fallback de \texttt{nivel}).}
\end{table}

% ===========================================================================
\section{Demais funções: cobertura de ramos e decisão}
% ===========================================================================

Para as outras funções aplicamos cobertura de \textbf{ramos/decisão} (cada
caminho do código exercitado ao menos uma vez). A tabela resume as decisões
cobertas e o número de casos.

\begin{table}[h]
\centering
\small
\renewcommand{\arraystretch}{1.15}
\begin{tabular}{@{}lp{8.6cm}c@{}}
\toprule
\textbf{Função} & \textbf{Ramos / decisões cobertos} & \textbf{Casos} \\
\midrule
\texttt{extractSubjectCodes}        & 1 código, vários códigos, normalização, sem código, vazio & 5 \\
\texttt{getCodigosEquivalentes}     & usa \texttt{expressao\_logica}; fallback p/ \texttt{expressao}; nenhum & 3 \\
\texttt{codigoContidoNaEquivalencia}& código vazio; com/sem \texttt{expressao\_logica}; contido/não & 4 \\
\texttt{getStatusPriority}          & 3 ramos: integralizada(3), MATR(2), demais(1) & 3 \\
\texttt{findSubjectMatch}           & 4 níveis de fallback (código/nome $\times$ obrig/opt) + sem match & 5 \\
\texttt{processMatchedDiscipline}   & nova; tipo optativa/obrig; duplicata prioridade maior/menor & 4 \\
\texttt{checkEquivalencies}         & satisfeita por lógica; por fallback; não integralizada; nenhuma & 4 \\
\texttt{mapEquivalenciasFromDb}     & mapeamento padrão; entrada nula; formato legado & 3 \\
\texttt{mapPreRequisitosFromDb}     & vazio; expansão por lógica; ramo \texttt{id\_materia\_requisito} & 3 \\
\bottomrule
\end{tabular}
\caption{Decisões cobertas pelos demais testes de unidade.}
\end{table}

% ===========================================================================
\section{Resultados de cobertura}
\label{sec:resultados}
% ===========================================================================

Medição com \texttt{jest --coverage} sobre
\texttt{src/controllers/fluxograma\_controller.ts}, antes e depois da nova suíte:

\begin{table}[h]
\centering
\small
\begin{tabular}{@{}lcccc@{}}
\toprule
\textbf{Métrica} & \textbf{Antes} & \textbf{Depois} & \textbf{Ganho} \\
\midrule
Statements & 10,08\% & 28,33\% & +18,3 p.p. \\
Branches   & 3,93\%  & 31,88\% & +28,0 p.p. \\
Functions  & 6,38\%  & 46,80\% & +40,4 p.p. \\
Lines      & 10,58\% & 28,23\% & +17,7 p.p. \\
\bottomrule
\end{tabular}
\caption{Cobertura do \texttt{fluxograma\_controller} antes/depois.}
\end{table}

A camada de funções auxiliares ficou praticamente toda coberta. O que resta sem
cobrir são os \emph{endpoints} HTTP assíncronos (linhas $\approx$415--1027), que
dependem fortemente do Supabase e exigem \emph{mocks} --- alvo de uma etapa
futura. A suíte completa do backend ficou em \textbf{84 testes / 7 suítes}, todos
passando.

\subsection*{Como reproduzir}

\begin{lstlisting}[language=bash]
cd no_fluxo_backend
npm ci
npx jest tests-ts/fluxograma_controller.whitebox.test.ts \
  --coverage --collectCoverageFrom='src/controllers/fluxograma_controller.ts'
\end{lstlisting}

\bigskip
\hrule
\bigskip
\noindent\small
Branch: \texttt{test/fluxograma-controller-coverage} \quad$\cdot$\quad
Arquivos: \texttt{src/controllers/fluxograma\_controller.ts} (export
\texttt{\_\_testing\_\_}) e \texttt{tests-ts/fluxograma\_controller.whitebox.test.ts}.

\end{document}
